\chapter{Motion generation and simulation-to-hardware transfer}
\label{ch:motion}

\Cref{ch:resultsmain} reports the headline numbers. This appendix gives the
mechanism.

\section{Why a per-joint clamp cannot stay on the line}
\label{sec:clamp-maths}

The limiter it replaced applied, per cycle and independently to each joint
$i$, $q_i^{(k+1)} = q_i^{(k)} + \operatorname{clip}(q_i^\star - q_i^{(k)},
-\Delta, +\Delta)$ with one step bound $\Delta$ shared by all seven joints.
Writing the required displacement as $\mathbf{e}=\mathbf{q}^\star-\mathbf{q}^{(k)}$,
staying on the straight joint-space line requires the applied step be
\emph{parallel} to $\mathbf{e}$; the clamp applies
$\operatorname{clip}(\mathbf{e},-\Delta,+\Delta)$ component-wise, which is
parallel to $\mathbf{e}$ only when every component is inside the bound or
every component saturates equally. In every other case the direction of
travel rotates and the end effector leaves the path. This is why sweeping
$\Delta$ over a factor of 35 does not drive the excursion to zero -- it
settles at \SI{68}{\milli\metre} -- and why the excursion is a property of
the algorithm, not its tuning.

\section{What replaced it}

Ruckig~\cite{berscheid2021ruckig}: time-optimal, jerk-limited, per-degree-of-freedom
limits, accepting a non-zero target velocity. Phase synchronisation drives
every joint from one scalar phase $\sigma(t)\in[0,1]$:
\begin{equation}
  \mathbf{q}(t) = \mathbf{q}^{(k)} + \sigma(t)\bigl(\mathbf{q}^\star-\mathbf{q}^{(k)}\bigr),
  \qquad \sigma(0)=0,\; \sigma(T)=1,
  \label{eq:phase-sync}
\end{equation}
so the commanded configuration is on the segment between endpoints by
construction for every $t$ -- unlike merely stretching each joint's own
profile to the slowest joint's duration (time synchronisation), which
measures \SI{9.5}{\milli\metre} of path deviation on the same case against
phase's \SI{0.02}{\milli\metre}. The acceleration and jerk limits are
\emph{assumed} and labelled so: \repo{joint\_limits.yaml} declares zero
acceleration limits for all fourteen joints, so there is nothing to read;
only the velocity limits are the robot's own.

Five defects were caught in the replacement, four of which would have
passed a conventional test suite: re-planning every cycle destroyed phase
synchronisation (a mid-profile re-plan cannot reproduce the profile it is
half through); floating-point noise forced a re-plan every cycle even after
fixing that, because the target was recomputed via a wrapped subtraction not
bit-identical to itself; the synchronised-clamp fallback reported zero
acceleration (because it does not model one) and so passed a continuity
check by fabricating the quantity being checked; the off-path metric had a
\SI{0.1}{\milli\metre} grid-resolution floor and was reporting the
instrument's resolution as the robot's behaviour; and an arrival test used
an absolute tolerance on travels spanning a factor of twenty-five.

\section{The simulation-to-hardware gap}
\label{sec:sim-to-real}

Of 252 recorded per-joint errors, 211 lie between \SI{0.25}{} and
\SI{0.30}{\degree} with the sign following the direction of travel -- a
\emph{bound}, not a spread, which identifies a terminal offset rather than
noise. Two candidate models, a terminal offset and a proportional gain
error, coincide at one travel distance and differ by a factor of five at
\SI{20}{\milli\metre}; propagating the recorded errors through the
manipulator Jacobian reproduces the measured Cartesian error to
\SI{90}{\percent} ($r=\num{0.923}$) under the offset model and not the gain
model. The data had been recorded for weeks in
\repo{recordings/baselines/arm\_directional\_calibration.json} before
anything read it.

An earlier report claimed the terminal error is speed-independent across a
fivefold range; the velocity limit was read once at bridge start-up and
seven of thirty-six runs moved \emph{faster} than commanded, so the three
speeds were three replicates of one condition and the speed dependence has
never actually been measured. The bridge deadband was \SI{1.0}{\degree}
against a measured \SI{0.305}{\degree} park -- a factor of 3.3 -- and is
\SI{0.10}{\degree} now, with an opt-in overshoot correction; neither has
been tried on hardware.

\section{The positioning error budget}

\begin{equation}
  \sigma_{\text{RSS}} = \sqrt{\textstyle\sum_i b_i^2}, \qquad
  \sigma_{\text{worst}} = \textstyle\sum_i b_i,
  \label{eq:budget}
\end{equation}

with $\sigma_{\text{worst}}$ the honest figure against the
\SI{30}{\milli\metre} capture gate, since the two largest terms are declared
constants rather than noise. Corrections -- the pad offset re-derived from
forward kinematics (\SI{13.45}{} to \SI{0.05}{\milli\metre}), the shell bias
removed via the measured support plane, and the terminal-offset overshoot
(joint term \SI{7.24}{} to \SI{1.88}{\milli\metre}) -- take the budget from
\SI{18.64}{} to \SI{12.91}{\milli\metre} RSS and \SI{33.39}{} to
\SI{19.99}{\milli\metre} worst case, \emph{provided} the camera-to-wrist
transform is small. \Cref{sec:hardware-session} measured it directly and it
is not.

\section{The reaction budget}

\begin{table}[htbp]
  \centering
  \caption[Reaction budget]{Limb travel during the measured wearer-tracking
  latency against the \SI{150}{\milli\metre} clearance floor. Speeds are
  literature ranges, labelled as assumptions; latencies are measured.}
  \label{tab:reaction}
  \begin{tabular}{@{}llrl@{}}
    \toprule
    Limb motion & Latency & Travel & Against the floor \\
    \midrule
    Deliberate reach & \SI{62}{\milli\second} & \SIrange{50}{93}{\milli\metre} & \SI{62}{\percent} of it \\
    Deliberate reach & \SI{562}{\milli\second} & \SIrange{450}{843}{\milli\metre} & exceeds it \\
    Quick reach & \SI{62}{\milli\second} & \SIrange{93}{155}{\milli\metre} & exceeds it \\
    Startle & \SI{62}{\milli\second} & \SIrange{155}{248}{\milli\metre} & exceeds it \\
    \bottomrule
  \end{tabular}
\end{table}

In five of six modelled limb motions the limb travels further than the
entire margin before the guard is aware anything has changed --
\cref{ch:discussionmain} treats this as the principal open safety question.
